% paper/sections/discussion.tex

\section{Discussion}
\label{sec:discussion}

\subsection{Why Classical Control Fails Under Fluid-Parameter Shift}

The $97$-pp \PID{} collapse is not a tuning failure---the gains were
optimised by exhaustive grid search on the training distribution.
It reflects a structural limitation of fixed-gain linear control:
proportional-integral terms cannot distinguish transient inertia-induced
error from persistent current-induced error.
Under test-distribution currents ($0.40$--$0.80$\,m/s), integral
windup saturates the thruster, causing overshoot; under high drag
($0.80$--$1.20$\,kg/m), the derivative term is underdamped, inducing
a limit cycle.
Any fixed-gain \PID{} tuned for nominal \AUV{} conditions will face
analogous failure wherever fluid parameters vary continuously.

\subsection{Why CDR Produces More Energy-Efficient Policies}
\label{sec:energy_mechanism}

The $1.13\%$ energy advantage of \CDR{} over \UDR{} follows directly
from their training structures.
\CDR{} begins in calm conditions where the reward penalty
$-0.02\|\mathbf{a}_t\|^2$ reinforces smooth, low-thrust trajectories;
as $\lambda$ grows, the agent extends this efficient strategy to
progressively harder conditions.
\UDR{} encounters high drag and strong currents from episode~1, where
near-maximum thrust is the only viable response; this aggressive habit
persists even when conditions do not require it.
The result: \CDR{} learns minimum viable thrust per condition;
\UDR{} learns a uniformly aggressive strategy that generalises
but never optimises.

\subsection{Deployment Implications}

A $1.13\%$ thrust reduction extends mission duration by the same
fraction for thrust-dominated \AUV{}s~\cite{fossen2011}---approximately
$2.7$\,min on a $4$-hour mission, enough for an additional survey pass.
We recommend that \AUV{} \RL{} benchmarks report energy efficiency
alongside success rate as a co-primary metric.

\subsection{Limitations}

\textbf{Simulation-only.}
Our fluid model follows~\citet{fossen2011} but omits turbulence and
wave-induced disturbances; hardware validation is left to future work.

\textbf{Limited seeds.}
With $n{=}3$, comparisons between closely matched conditions have
limited power; future work with $n{\geq}10$ seeds would strengthen
evidence.

\textbf{Shared expansion window.}
\CDR{} expands all six parameters simultaneously; per-parameter
expansion~\cite{akkaya2019} may provide finer curriculum control at
the cost of additional tuning complexity.